\uuid{vVeU}
\titre{Matrice d'une application linéaire/Matrice de passage/Espace des fonctions}
\auteur{hanani ; mzari}
\organisation{exo7}
\licence{CC BY-NC-SA 4.0}
\datecreate{2020-01-20}
\niveau{L1}
\module{Algèbre}
\chapitre{Matrice}
\sousChapitre{Matrice et application linéaire}
\difficulte{3}
\theme{Matrice d'une application linéaire, Matrice de passage, Espace des fonctions}
\exo7qcmid{lille-2.matrice-application.017}

\contenu{
\begin{question}
\qclassification{108.03}
\qtags{motcle=Matrice d'une application linéaire/Matrice de passage/Espace des fonctions}
Soit ${\cal F}$ l'espace vectoriel des fonctions réelles engendré par les fonctions $f_1, \; f_2$ et  $f_3$ définies par  :
$f_1(x)=1, \; f_2(x)=\cos x$ et $f_3(x)= \sin x$. On munira ${\cal F}$ des bases ${\cal B}=\{f_1,f_2,f_3\}$ et ${\cal B}'=\{f_1,f_2+f_3,f_2-f_3\}$. On notera $P$ la matrice de passage de la base $\cal {B}$ à la base $\cal {B}'$ et $Q$ la matrice de passage de la base ${\cal {B'}}$ à la base ${\cal {B}}$.
\vskip1mm
Définition : Soit $E$ un espace vectoriel muni de deux bases ${\cal {B}}$ et ${\cal {B}}'$. La matrice de passage de la base ${\cal {B}}$ à la base  ${\cal {B}}'$ est la matrice de l'identité de $E$ de la base ${\cal {B}}'$ à la base  ${\cal {B}}$. Autrement dit, c'est la matrice dont la jième colonne est constituée des coordonnées du jième vecteur de la base ${\cal {B}'}$ dans la base  ${\cal {B}}$.
\vskip1mm
Quelles sont les assertions vraies ?
\begin{answers}
\good{$
P =\left(\begin{array}{rcc}
1&0&0\\
0&1&1\\
0&1&-1\\
\end{array}\right).$}
\bad{$
P = \frac{1}{2}\left(\begin{array}{rcc}
2&0&0\\
0&1&1\\
0&1&-1\\
\end{array}\right).$}
\good{$Q= \frac{1}{2}\left(\begin{array}{rcc}
2&0&0\\
0&1&1\\
0&1&-1\\
\end{array}\right).$}
\bad{La matrice de l'application identité de ${{\cal F}}$ de la base ${\cal {B}}$ à la base ${\cal {B'}}$ est :
$$\left(\begin{array}{rcc}
1&0&0\\
0&1&1\\
0&1&-1\end{array}\right).$$}
\end{answers}
\vskip2mm
\begin{explanations} Par d\'efinition, $P$ est la matrice de l'application identité de ${{\cal F}}$ de la base ${\cal {B}'}$ à la base ${\cal {B}}$ et $Q$ est la matrice de l'application identité de ${{\cal F}}$ de la base ${\cal {B}}$ à la base ${\cal {B'}}$. Donc
$$P= \left(\begin{array}{rcc}
1&0&0\\
0&1&1\\
0&1&-1\\
\end{array}\right)\quad \mbox{donc}\quad Q= \frac{1}{2}\left(\begin{array}{rcc}
2&0&0\\
0&1&1\\
0&1&-1\end{array}\right).$$
\end{explanations}
\end{question}
}
